% ============================================================================
% CAPÍTULO 3: DELIMITACIÓN DEL OBJETO DE ESTUDIO
% ============================================================================
%
% REQUISITOS SEGÚN LA FACULTAD:
%
% PLANTEAMIENTO DEL PROBLEMA Y PREGUNTA DE INVESTIGACIÓN:
% • Describir el problema que se observa en su contexto y fundamentarlo con
%   datos estadísticos
% • Expresa la relación entre dos o más conceptos o variables de forma clara
%   y precisa a través de una aseveración o interrogante congruente con el
%   enfoque de la investigación y el objetivo de investigación
%
% HIPÓTESIS:
% • Contiene una predicción o explicación provisional de una relación entre
%   dos o más variables la cual es susceptible de someterse a prueba
% • Enunciados teóricos supuestos, no verificados, pero probables referentes
%   a variables o a relación entre variables
% • Soluciones probables previamente seleccionadas, al problema planteado,
%   que se proponen para que, durante el proceso de la investigación, sean
%   confirmadas o no por los hechos
%
% OBJETIVOS:
% Objetivo General:
% • Señala el alcance de la investigación, acorde a la pregunta de investigación,
%   al enfoque y método de investigación
%
% Objetivos Específicos:
% • Contempla lo que se pretende obtener y qué hacer con los resultados
% • Se trata de hacer explícito el para qué de la investigación
% ============================================================================

\chapter{Delimitación del Objeto de Estudio}
\label{cap:delimitacion}

% ----------------------------------------------------------------------------
% SECCIÓN 3.1: Planteamiento del Problema
% ----------------------------------------------------------------------------
\section{Planteamiento del Problema}
\label{sec:planteamiento_problema}

% INSTRUCCIONES:
% - Describe el problema observado EN SU CONTEXTO
% - FUNDAMENTA CON DATOS ESTADÍSTICOS (prevalencias, incidencias, etc.)
% - Sé específico sobre el contexto (local, nacional, internacional)
% - Conecta con lo presentado en el marco teórico

\subsection{El Problema en Contexto}
\label{subsec:problema_contexto}

[Describe el problema específico que abordas]

\textbf{Contexto Internacional:}

A nivel mundial, \textit{[estadística/dato]} \cite{fuente_internacional}. Esto representa \textit{[magnitud del problema]}.

\textbf{Contexto Nacional (México):}

En México, \textit{[estadística nacional]} según \cite{fuente_nacional_ej_INEGI}. Específicamente, \textit{[detalle relevante]}.

\textbf{Contexto Local (Chihuahua):}

En el estado de Chihuahua, \textit{[estadística local]} \cite{fuente_local}. [Describe la situación específica en tu contexto].

\subsection{Delimitación Específica del Problema}
\label{subsec:delimitacion_problema}

El problema específico que se aborda en esta investigación es:

\textit{[Describe de forma precisa y delimitada el problema. Debe quedar claro qué aspecto específico abordarás, en qué población, en qué contexto, etc.]}

Este problema se manifiesta en \textit{[consecuencias observables]} y afecta a \textit{[población específica]}.

% ----------------------------------------------------------------------------
% SECCIÓN 3.2: Pregunta de Investigación
% ----------------------------------------------------------------------------
\section{Pregunta de Investigación}
\label{sec:pregunta_investigacion}

% INSTRUCCIONES:
% - Debe expresar la relación entre DOS O MÁS conceptos/variables
% - Debe ser clara y precisa
% - Debe ser congruente con el enfoque de investigación
% - Debe ser RESPONDIBLE con tu metodología

Con base en el problema planteado, la pregunta que guía esta investigación es:

\begin{center}
\fbox{%
\parbox{0.9\textwidth}{%
\centering
\textbf{¿[Tu pregunta de investigación]?}
}%
}
\end{center}

\textbf{Ejemplo de estructura:}
\begin{itemize}
    \item \textit{¿Cuál es la relación entre [Variable 1] y [Variable 2] en [Población] durante [Tiempo/Contexto]?}
    \item \textit{¿Cómo influye [Variable Independiente] en [Variable Dependiente] en [Contexto específico]?}
    \item \textit{¿Qué efectos tiene [Intervención] sobre [Variable de resultado] en [Población objetivo]?}
\end{itemize}

% ----------------------------------------------------------------------------
% SECCIÓN 3.3: Hipótesis
% ----------------------------------------------------------------------------
\section{Hipótesis}
\label{sec:hipotesis}

% INSTRUCCIONES:
% - Contiene una PREDICCIÓN o explicación provisional
% - Debe ser susceptible de SOMETERSE A PRUEBA
% - Especifica la relación esperada entre variables
% - Debe ser congruente con tu marco teórico

\subsection{Hipótesis de Investigación (Hi)}
\label{subsec:hipotesis_investigacion}

\textbf{Hi:} \textit{[Enunciado que predice la relación entre variables]}

\textbf{Ejemplo de estructura:}
\begin{itemize}
    \item \textit{Existe una relación significativa entre [Variable 1] y [Variable 2] en [Población]}
    \item \textit{[Variable Independiente] tiene un efecto positivo/negativo sobre [Variable Dependiente]}
    \item \textit{[Grupo experimental] presentará mejores resultados en [Variable] comparado con [Grupo control]}
\end{itemize}

\subsection{Hipótesis Nula (Ho)}
\label{subsec:hipotesis_nula}

\textbf{Ho:} \textit{[Enunciado que niega la relación o afirma la ausencia de efecto]}

\textbf{Ejemplo:}
\begin{itemize}
    \item \textit{No existe relación significativa entre [Variable 1] y [Variable 2]}
    \item \textit{[Variable Independiente] no tiene efecto sobre [Variable Dependiente]}
\end{itemize}

\subsection{Hipótesis Alternativas (Ha) [Si aplica]}
\label{subsec:hipotesis_alternativas}

Si tienes múltiples hipótesis específicas, enuméralas:

\begin{enumerate}
    \item \textbf{Ha1:} \textit{[Primera hipótesis alternativa]}
    \item \textbf{Ha2:} \textit{[Segunda hipótesis alternativa]}
\end{enumerate}

% ----------------------------------------------------------------------------
% SECCIÓN 3.4: Objetivos
% ----------------------------------------------------------------------------
\section{Objetivos}
\label{sec:objetivos}

% INSTRUCCIONES:
% - Objetivo General: alcance de la investigación, acorde a pregunta e hipótesis
% - Objetivos Específicos: qué pretendes obtener, para qué
% - Usa verbos en INFINITIVO: Determinar, Evaluar, Analizar, Comparar, etc.
% - Deben ser MEDIBLES y ALCANZABLES

\subsection{Objetivo General}
\label{subsec:objetivo_general}

\textbf{Verbo en infinitivo} \textit{[la relación/efecto/asociación]} entre \textit{[variables]} en \textit{[población]} mediante \textit{[metodología general]} para \textit{[finalidad]}.

\textbf{Ejemplos de verbos para objetivo general:}
\begin{itemize}
    \item Determinar, Evaluar, Analizar, Identificar, Comparar, Establecer, Describir, Examinar, Valorar
\end{itemize}

\textbf{Ejemplo completo:}

\textit{Determinar la relación entre el nivel de actividad física y el comportamiento sedentario en adultos mayores de Chihuahua mediante un estudio transversal para proponer estrategias de intervención.}

\subsection{Objetivos Específicos}
\label{subsec:objetivos_especificos}

% INSTRUCCIONES:
% - Deben ser pasos específicos para alcanzar el objetivo general
% - Generalmente 3-5 objetivos específicos
% - Deben seguir un orden lógico (metodológico o temporal)

\begin{enumerate}[label=\textbf{OE\arabic*:}]
    \item \textbf{[Verbo en infinitivo]} \textit{[qué harás específicamente]} mediante \textit{[cómo lo harás]} para \textit{[para qué]}.
    
    \textbf{Ejemplo:} \textit{Caracterizar el perfil sociodemográfico de la población de estudio mediante estadística descriptiva para identificar factores asociados.}
    
    \item \textbf{[Verbo]} \textit{[acción específica 2]}...
    
    \item \textbf{[Verbo]} \textit{[acción específica 3]}...
    
    \item \textbf{[Verbo]} \textit{[acción específica 4]}...
\end{enumerate}

\textbf{Sugerencia de estructura lógica:}
\begin{itemize}
    \item OE1: Caracterizar/Describir la población
    \item OE2: Medir/Evaluar variables principales
    \item OE3: Analizar relaciones/asociaciones
    \item OE4: Proponer/Generar conclusiones o recomendaciones
\end{itemize}

% ----------------------------------------------------------------------------
% RESUMEN DE CONGRUENCIA
% ----------------------------------------------------------------------------
\section{Congruencia entre Problema, Pregunta, Hipótesis y Objetivos}
\label{sec:congruencia}

% INSTRUCCIONES:
% - Verifica que todo esté alineado
% - Esta sección ayuda a visualizar la coherencia de tu investigación

La \Cref{tab:congruencia} muestra la congruencia entre los elementos fundamentales de esta investigación:

\begin{table}[htbp]
    \centering
    \caption{Congruencia entre elementos de la investigación}
    \label{tab:congruencia}
    \small
    \begin{tabular}{@{}lp{10cm}@{}}
        \toprule
        \textbf{Elemento} & \textbf{Enunciado} \\
        \midrule
        
        \textbf{Problema} & 
        [Síntesis del problema en 1-2 líneas] \\
        
        \midrule
        
        \textbf{Pregunta} & 
        [Tu pregunta de investigación] \\
        
        \midrule
        
        \textbf{Hipótesis (Hi)} & 
        [Tu hipótesis de investigación] \\
        
        \midrule
        
        \textbf{Objetivo General} & 
        [Tu objetivo general] \\
        
        \midrule
        
        \textbf{Objetivos Específicos} & 
        OE1: [Resumen]\\
        OE2: [Resumen]\\
        OE3: [Resumen]\\
        OE4: [Resumen] \\
        
        \bottomrule
    \end{tabular}
\end{table}

Como se observa en la \Cref{tab:congruencia}, existe coherencia entre el problema identificado, la pregunta de investigación formulada, la hipótesis propuesta y los objetivos planteados, lo cual garantiza la viabilidad y pertinencia de este estudio.

% ============================================================================
% NOTAS FINALES:
% ============================================================================
%
% RECUERDA:
% 1. Problema: fundamentado con DATOS ESTADÍSTICOS
% 2. Pregunta: Clara, precisa, respondible
% 3. Hipótesis: Predicción susceptible de prueba
% 4. Objetivo General: Alcance de la investigación
% 5. Objetivos Específicos: Pasos medibles y alcanzables
% 6. TODO debe estar ALINEADO y ser CONGRUENTE
%
% ============================================================================
